\documentclass[10pt]{article}

\usepackage[utf8]{inputenc}

\usepackage[T1]{fontenc}

\usepackage{amsmath}

\usepackage{amsfonts}

\usepackage{amssymb}

\usepackage[version=4]{mhchem}

\usepackage{stmaryrd}

\usepackage{graphicx}

\usepackage[export]{adjustbox}

\graphicspath{ {./images/} }



\begin{document}

БОРА - ПАЙЕРЛСА - ПЛАЧЕКА СООТНОШЕНИЕ см. Оптическая теорема.



БОРЕЛЕВСКОЕ ПРЕОБРАЗОВАНИЕ - см. Равномерно пригодные приближения в квантовой механике.



БОРЕЛЯ МЕТОД - метод суммирования асимптотических рядов. См. Равномерно пригодные приближения в квантовой механике.



БOPHА - ГРИНА - ИВОНА УРАВНЕНИЕ - приближенное замкнутое нелинейное уравнение для парной функции распределения (см. Парная корреляционная функция) классической равновесной системы. Система представляет собой ансамбль $N$ тождественных точечных частиц массой $m$ каждая, находящихся в ящике объема $V$. Гамильтониан системы имеет вид



\[

H=\sum_{j=1}^{N} \frac{p_{j}^{2}}{2 m}+\sum \sum_{1 \leqslant j<n \leqslant N} v\left(r_{j n}\right),

\]



где $r_{j n}=\left|\overrightarrow{q_{j}}-\overrightarrow{q_{n}}\right|, \overrightarrow{q_{1}}, \ldots, \overrightarrow{q_{N}}, \overrightarrow{p_{1}}, \ldots, \overrightarrow{p_{N}}-$ координаты и импульсы частиц системы. Б.- Г.- И. у. получается из цепочки уравнений для частичных функций распределения после применения феноменологич. суперпозиционного приближения Кирквуда, когда трехчастичную конфигурационную функцию распределения выражают через двухчастичную



\[

n_{3}\left(\overrightarrow{q_{1}}, \overrightarrow{q_{2}}, \overrightarrow{q_{3}}\right)=n_{2}\left(\overrightarrow{q_{1}}, \overrightarrow{q_{3}}\right) n_{2}\left(\overrightarrow{q_{2}}, \overrightarrow{q_{3}}\right) n_{2}\left(\overrightarrow{q_{1}}, \overrightarrow{q_{2}}\right)

\]



Б. - Г.- И. у. имеет вид:



\[

\begin{aligned}

& -\theta \vec{\nabla}_{1}\left[\ln n_{2}\left(r_{12}\right)\right]=\vec{\nabla}_{1} v\left(r_{12}\right)+ \\

& +n \int d \vec{q}_{3}\left[\vec{\nabla}_{1} v\left(r_{13}\right)\right] n_{2}\left(r_{13}\right) n_{2}\left(r_{23}\right),

\end{aligned}

\]



где $\theta$ - температура, $n$ - средняя плотность числа частиц, $\vec{\nabla}_{1} \equiv \partial / \partial \overrightarrow{q_{1}}$. Б.- Г.- И. у. впервые приведено в [1]..



Jum.: [1] B orn M., Gre e n H.S., «Proc. Roy. Soc.», A 188, 1946, v. 10; [2] Бале с ку Р., Равновесная и неравновесная статистическая механика, пер. с англ., т. 1, М., 1978. БОРНА - ИНФЕЛЬДА УРАВНЕНИЯ - Максвела уравнения для поля точечных источников в «вакууме» с нулевой проводимостью и материальными уравнениями



\[

\boldsymbol{D}=\frac{\boldsymbol{E}-G \boldsymbol{B}}{L}, \quad \boldsymbol{H}=\frac{\boldsymbol{B}-G \boldsymbol{E}}{L},

\]



где $L=\sqrt{1+F-G^{2}}-$ лагранжиан Б.- И.у., $F=\frac{B^{2}-E^{2}}{b^{2}}$, $G=\frac{B E}{b^{2}}$ и $b-$ константа порядка $10^{10} \mathrm{k} / \mathrm{M}^{2}$. Б.- И. у.- модель электродинамики, в к-рой источники отождествляются с сингулярными точками $D$, так что траектории источников определяются из уравнений порождаемого ими поля. При этом энергия, импульс и спин источников принимаются равными соответствующим величинам его поля.



Jum.: [1] B orn M., Infeld L., «Proc. Roy. Soc. Ser. A», 1934 v. 144, № 852 , p. $425-52$; v. 147, № 862, p. 522-47; 1935, v. 150 , № 869, p. 141-66; [2] Pryc e M. H. L., там же, 1935, v. 150, № 869,



\includegraphics[max width=\textwidth, center]{2023_07_08_c57a17bf2025ab8ed6c6g-1}

БOPHА - OПกЕНТЕЙМЕРА МЕТОД - метод, вПервЫе сформулированный в [1] и являющийся основным методом описания и классификации молекулярных спектров, медленных атомно-молекулярных столкновений. Используется также в динамич. теории кристаллич. решеток. Представляет собой один из вариантов адиабатической теории возмущений и основан на выделении в квантовой системе «быстрой» (электроны) и «медленной» (ядра) подсистем (см. $A \partial u$ абатическое приближение). Такое разделение становится возможным вследствие большой разницы в массах электронов и ядер. В результате исходная задача решается в два этапа: вначале определяются состояния электронов при фиксированных координатах ядер (электронные термы), а затем учитывается движение последних. Пусть имеется квантовая система, состоящая из нек-рого числа электронов с массой $m$ и атомных ядер, в общем случае с различными массами $M_{\alpha}$. Совокупность координат всех электронов относительно центра инерции системы $\boldsymbol{r} \equiv\left\{\boldsymbol{r}_{1}, \boldsymbol{r}_{2}, \ldots, \boldsymbol{r}_{i}, \ldots\right\}$, а $\boldsymbol{R} \equiv\left\{\boldsymbol{R}_{1}, \boldsymbol{R}_{2}, \ldots, \boldsymbol{R}_{\alpha}, \ldots\right\}$ - совокупность координат ядер. Гамильтониан такой системы имеет вид



\[

\hat{H}=\hat{T}_{r}+\hat{T}_{R}+\hat{V}(\boldsymbol{r}, \boldsymbol{R}),

\]



где



\[

\hat{T}_{r}=-\frac{\hbar^{2}}{2 m} \sum_{i} \frac{\partial^{2}}{\partial r_{i}^{2}} ; \quad \hat{T}_{R}=-\hbar^{2} \sum_{\alpha} \frac{1}{M_{\alpha}} \frac{\partial^{2}}{\partial R_{\alpha}^{2}}

\]



\begin{itemize}

  \item соответственно операторы кинетич. энергии электронов и ядер, а $V(\boldsymbol{r}, \boldsymbol{R})$ - оператор потенциальной энергии взаимодействия между частицами.

\end{itemize}



Поскольку в реальных системах выполнено условие $M_{\alpha} \gg m$, то оператор кинетич. энергии $\hat{T}_{R}$ тяжелых частиц (ядер) можно рассматривать как малое возмущение, что и отличает Б. - О. м. от более распространенных схем квантовомеханич. теории возмущений, в К-рых малой предполагается потенциальная энергия взаимодействия. Рассматриваемое возмущение имеет сингулярный характер и, как следствие, асимптотич. характер связанных с ним разложений, поскольку малый параметр входит в ту часть гамильтониана, к-рая содержит старшие производные по группе переменных системы.



В нулевом порядке Б.- О. м. задача определения стационарных состояний системы сводится к решению уравнения Шредингера для волновых функций электронов



\[

\left[\hat{T}_{r}+V(\boldsymbol{r}, \boldsymbol{R})\right] \varphi_{n}(\boldsymbol{r}, \boldsymbol{R})=\varepsilon_{n}(\boldsymbol{R}) \varphi_{n}(\boldsymbol{r}, \boldsymbol{R}),

\]



в к-ром координаты ядер являются параметрами, входящими как в волновые функции, так и в энергетич. уровни $\varepsilon_{n}(\boldsymbol{R})$, определяющие электронный терм рассматриваемой системы. Здесь индекс $n$ определяет совокупность квантовых чисел, соответствующих стационарному состоянию системы электронов.



Если предположить, что полный набор решений уравнения (2) при произвольных значениях $\boldsymbol{R}$ известен, то собственные функции полного гамильтониана (1), то есть решения уравнения



\[

(\hat{H}-E) \Psi(\boldsymbol{r}, \boldsymbol{R})=0

\]



можно искать в виде разложения по функциям $\varphi_{n}(\boldsymbol{r}, \boldsymbol{R})$ :



\[

\Psi(\boldsymbol{r}, \boldsymbol{R})=\sum_{n} \Phi_{n}(\boldsymbol{R}) \varphi_{n}(\boldsymbol{r}, \boldsymbol{R}),

\]



причем в общем случае суммирование по $n$ включает в себя как дискретный, так и непрерывный спектры «быстрой» подсистемы.



Подстановка указанного разложения в уравнение (3) приводит к следующей системе уравнений для коэффициентных функций $\Phi_{n}(\boldsymbol{R})$ (см. [2]):



\[

\left[\hat{T}_{R}+\varepsilon_{n}(R)+\hat{\Omega}_{n n}(R)-E\right] \Phi_{n}(R)=\sum_{m \neq n} \hat{\Omega}_{m n}(R) \Phi_{m}(R)

\]



где матричный оператор



\[

\begin{aligned}

& \hat{\Omega}_{m n}(\boldsymbol{R})=\sum_{\alpha} \frac{\hbar^{2}}{M_{\alpha}} \int \varphi_{m}^{*}(\boldsymbol{r}, \boldsymbol{R}) \frac{\partial}{\partial \boldsymbol{R}_{\alpha}} \varphi_{n}(\boldsymbol{r}, \boldsymbol{R}) d \boldsymbol{r} \frac{\partial \boldsymbol{R}}{\partial \boldsymbol{R}_{\alpha}}- \\

&-\int \varphi_{m}^{*}(\boldsymbol{r}, \boldsymbol{R}) \hat{T}_{R} \varphi_{n}(\boldsymbol{r}, \boldsymbol{R}) d \boldsymbol{r} .

\end{aligned}

\]



В точных уравнениях (4) адиабатич. разложение Б.- О. м. для построения состояний «медленной» подсистемы определяется рядом по степеням оператора $\hat{\Omega}$, причем в первом порядке система (4) распадается на независимые уравнения



\[

\left[\hat{T}_{R}+\varepsilon_{n}(R)+\hat{\Omega}_{n n}(R)-E_{n v}^{(0)}\right]_{n N}^{(0)}(R)=0

\]



для каждого электронного терма. Здесь $v-$ совокупность квантовых чисел, определяющих энергетич. уровни системы ядер в поле



\[

\hat{U}_{n}(\boldsymbol{R})=\varepsilon_{n}(\boldsymbol{R})+\hat{\Omega}_{m n}(\boldsymbol{R})

\]





\end{document}